\documentclass[journal]{IEEEtran}
\usepackage[T1]{fontenc}
\usepackage{graphicx}
\usepackage{booktabs}
\usepackage{amsmath}
\usepackage{array}
\usepackage{longtable}
\usepackage{tabularx}
\usepackage{enumitem}
\usepackage[hidelinks]{hyperref}
\usepackage{url}
\graphicspath{{../figures/}}
\newcolumntype{Y}{>{\raggedright\arraybackslash}X}
\begin{document}
\title{Multi-View and Multimodal Perception for Class-Wise Fruit Inventories: A Design-Space Review}
\author{Muhammad~Zainal~Muttaqin and Fatma~Indriani\thanks{Department of Computer Science, Universitas Lambung Mangkurat, Banjarbaru, Indonesia.}\thanks{Search, prioritization, and evidence-ledger audit executed 10 August 2026.}}
\maketitle
\begin{abstract}
Turning repeated observations of one tree into a single inventory of unique fruits grouped by class requires more than per-image detection. This design-space review treats oil-palm fresh fruit bunches as the principal agricultural case and transfers identity mechanisms from fruit, orchard, and non-agricultural multi-view perception. Scopus and OpenAlex searches used seven reproducible query families, explicit years, inclusion and exclusion rules, DOI or title-year deduplication, and title-abstract triage. The search produced 21,066 resolved master records; 20,035 advanced after abstract screening. These records were ranked by deterministic title-abstract signals and diversified into a 250-record shortlist with a 60-record first wave, rather than being read one by one. At the reporting cutoff, 44 targeted full-text studies supplied complete evidence rows. The synthesis distinguishes intra-view detection, statistical correction, appearance matching, temporal tracking, geometric association, and learned multi-view association. It generalizes class attributes beyond maturity to size, quality, disease, cultivar, and harvest readiness, while retaining the requirement that attributes attach to persistent instances. The resulting contribution is a testable design space for duplicate-aware, class-wise fruit inventories and an explicit boundary between search candidates and verified evidence.
\end{abstract}
\begin{IEEEkeywords}
multi-view perception, multimodal fusion, fruit counting, instance identity, oil palm, design-space review
\end{IEEEkeywords}
% Final revision candidate based on the eight lecturer revision points.
% Search candidates and verified full-text evidence are kept as separate registers.

\section{Introduction}
\label{sec:intro2}

An agricultural vision system may observe the same tree from several positions, frames, or sensors. The operational output is not the number of detections emitted by those images. It is an inventory of unique fruit instances, grouped by an attribute class. For a tree with instance set $I$, the desired count for class $c$ is the number of instances whose assigned attribute is $c$. A detection that is missed, merged with a neighbour, assigned the wrong class, or counted again from another view produces a different error. These errors cannot be diagnosed from image-level accuracy alone.

This review addresses the task as a design problem:
\begin{quote}
\emph{How can many observations of one tree be converted into a class-wise inventory in which each fruit or fruit bunch is counted once?}
\end{quote}
The principal agricultural case is the oil-palm fresh fruit bunch (FFB). The class attribute is not restricted to maturity. Maturity is one example; size, quality, disease, defect, cultivar, or harvest readiness are also possible attributes. All of them remain instance-dependent because a class-wise count is undefined until an individual fruit or bunch has been localized and assigned an identity.

The scope is agricultural fruit perception, but the mechanism search is deliberately wider. Tracking, person or vehicle re-identification, multi-camera association, visual odometry, structure from motion (SfM), and 3D object detection were developed largely outside orchards. They are included as transferable evidence when they answer a mechanism question about identity, geometry, or duplicate rejection. Their results are not presented as agricultural performance.

The review makes four contributions. First, it reports a reproducible Scopus and OpenAlex search with query strings, date bounds, selection rules, and a deterministic deduplication record. Second, it organizes the task around six identity mechanisms and the acquisition assumptions each mechanism requires. Third, it positions previous oil-palm, fruit-counting, and fruit-ripeness reviews against the cross-view identity problem. Fourth, it provides a row-level synthesis matrix in which method, modality, evaluation context, identity mechanism, class attribute, finding, and transfer boundary are visible together.

The paper does not propose one architecture as the universal solution. It asks when a simple statistical correction, temporal tracking, appearance matching, geometric reconstruction, or learned multi-view model is justified by the acquisition regime. This distinction is important for oil palm because the same tree can be observed from four or eight positions while sensor calibration, depth quality, illumination, and occlusion vary at the same time.

\section{Review questions and reproducible protocol}
\label{sec:protocol2}

\subsection{Review type and questions}

This work is a structured design-space review reported with a PRISMA-style search and selection trace. It is not a meta-analysis. The included tasks, object definitions, sensor modalities, and evaluation metrics are too heterogeneous for a pooled effect size. The protocol was not registered before execution. The review questions are:

\begin{enumerate}[leftmargin=*,itemsep=1pt]
  \item What operations are required to turn repeated observations into a unique, class-wise inventory?
  \item Which identity mechanisms have been used for fruit or for transferable non-agricultural objects?
  \item Which assumptions about overlap, motion, calibration, appearance, and view order make each mechanism plausible?
  \item What evidence exists for oil-palm and other tree-fruit systems, and what remains a transfer hypothesis?
  \item Which measurements would distinguish a better detector from a better inventory system?
\end{enumerate}

\subsection{Sources, date range, and query families}

Scopus was used as the subscription index available in the captured institutional session. OpenAlex was used as an open complementary index. Web of Science was not available in the recorded run; this deviation is stated rather than hidden. Google Scholar was reserved for possible backward or forward snowballing and was not used as a primary count source because its result export is not stable. The date window was 2015--2026. In Scopus the upper and lower bounds were written explicitly as \texttt{PUBYEAR > 2014 AND PUBYEAR < 2027}; the equivalent Web of Science syntax is recorded in the protocol for a future rerun.

Seven query families were used. Q1 targets fruit inventories from multiple observations. Q2 targets cross-view identity and duplicate rejection, including non-agricultural objects. Q3 targets plant geometry and multimodal sensing. Q4 targets class attributes attached to instances. Q5 targets previous reviews for positioning. Q6 targets oil-palm vision and machine learning. Q7 targets single-view fruit counting and yield estimation, which Q1 would otherwise miss. The exact strings used in the Scopus run are reproduced in Appendix~\ref{app:queries2}. The canonical query file is \path{literature/search/scopus-queries.md}; its count export and SHA-256 fingerprints are in \path{literature/search/scopus-counts-2026-08-09.csv}.

\begin{table}[t]
\centering
\caption{Search yields before cross-source deduplication.}
\label{tab:searchcounts2}
\small
\begin{tabular}{lrrrr}
\toprule
Query & Scopus & OpenAlex & Main focus & Known item\\
\midrule
Q1 & 2,001 & 1,849 & multi-view fruit inventory & SawitMVC\\
Q2 & 2,789 & 1,991 & identity and duplicates & mechanism set\\
Q3 & 5,648 & 6,422 & plant geometry and sensors & Gené-Molá\\
Q4 & 2,422 & 2,757 & per-instance attributes & class labels\\
Q5 & 662 & 1,143 & previous reviews & review set\\
Q6 & 995 & 1,177 & oil-palm vision & FFB studies\\
Q7 & 1,205 & 1,317 & single-view fruit load & MangoYOLO\\
\midrule
Total exported & 15,722 & 16,656 & & \\
\bottomrule
\end{tabular}
\end{table}

The Scopus known-item checks passed for SawitMVC in Q1, Gené-Molá et al. in Q3, and MangoYOLO in Q7. The checks demonstrate that the subject-area filter did not silently remove the three anchors. They do not demonstrate that either index is complete.

\subsection{Eligibility and elimination rules}

The following rules were applied to title, abstract, and full text in sequence. A source was eligible for the focused evidence register when all five inclusion conditions were met:

\begin{description}[leftmargin=*,style=nextline,itemsep=2pt]
  \item[IC1.] It is a peer-reviewed research article, conference paper, or openly available preprint with a verifiable full text and a quantitative result.
  \item[IC2.] Its visual output is per instance, such as object detection, instance segmentation, tracking, localization, or cross-observation association.
  \item[IC3.] Its publication year is 2015--2026, except for foundational Register B citations used only for context.
  \item[IC4.] It is written in English or Indonesian.
  \item[IC5.] The full text identifies a dataset or acquisition protocol and an evaluation metric or quantitative comparison.
\end{description}

The elimination rules were EC1, output only at image or scene level without per-instance identity; EC2, no quantitative evaluation; EC3, full text unavailable; EC4, duplicate or version of an already retained record; EC5, editorial, abstract-only, poster, patent, or non-research document; and EC6, outside the date range. A title or abstract was excluded automatically only when the signal was high confidence. Ambiguous records were retained for full-text assessment.

\subsection{Deduplication and screening trace}

The raw files were kept unchanged. Within each query, exact duplicate identifiers were removed. Across sources, the primary key was a normalized DOI. Records without a DOI used normalized title plus year. A title-year collision was not merged when author or venue evidence conflicted. The resolver then merged only groups with strong author overlap or an explicit DOI and venue agreement; conservative groups remained separate when evidence was insufficient.

The seven Scopus exports contained 15,722 rows and the seven OpenAlex exports contained 16,656 rows. Together they contained 32,378 raw rows, 32,222 rows after within-query deduplication, and 22,269 records after exact-key union. The resolver processed 1,030 merge components and reduced the working master to 21,066 records. The complete decision records are stored in the following files:
\begin{quote}
\path{literature/search/derived/dedup-resolution-2026-08-09.csv}\par
\path{literature/search/derived/master-screening-2026-08-09.csv}
\end{quote}

High-confidence title screening removed 242 EC5 records and one EC6 record, leaving 20,823 records for abstract screening. Abstract screening removed 640 EC1 records and 148 EC5 records, leaving 20,035 search-pool candidates. The 20,035 records were not treated as a requirement to read 20,035 papers one by one. A deterministic title-abstract score ranked core identity and inventory signals, fruit and oil-palm targets, instance perception, geometry, tracking, and prior reviews, while penalizing global-only outputs and non-target domains. A diversified shortlist of 250 records was selected, with a first targeted wave of 60. By the manuscript cutoff, 44 records had complete full-text evidence in the ledger and 16 had been excluded after review.

\begin{figure}[t]
\centering
\small
\begin{tabular}{c}
\fbox{\parbox{0.78\columnwidth}{\centering 32,378 raw records\\Scopus 15,722 + OpenAlex 16,656}}\\[2pt]
$\downarrow$\\[-1pt]
\fbox{\parbox{0.78\columnwidth}{\centering 22,269 exact-key master records}}\\[2pt]
$\downarrow$\\[-1pt]
\fbox{\parbox{0.78\columnwidth}{\centering 21,066 resolved working-master records}}\\[2pt]
$\downarrow$\\[-1pt]
\fbox{\parbox{0.78\columnwidth}{\centering 20,035 abstract-advanced search-pool records}}\\[2pt]
$\downarrow$\\[-1pt]
\fbox{\parbox{0.78\columnwidth}{\centering 250 targeted shortlist; 60 first wave; 44 included evidence studies}}
\end{tabular}
\caption{Search, prioritization, and targeted full-text trace. The 20,035 record pool is ranked and filtered before full-text review; it is not a promise of 20,035 individual paper reviews.}
\label{fig:flow3}
\end{figure}

\subsection{Prioritization before full-text review}

After abstract screening, the 20,035 records were treated as a search pool rather than as a mandatory reading list. The ranking script \path{tools/build\_literature\_priority\_shortlist.py} scored reproducible title and abstract signals. Positive signals covered unique inventory, duplicate resolution, re-identification, cross-view association, tracking, SfM or MVS, point clouds, RGB-D, depth, 3D reconstruction, fruit, FFB, and oil palm. Penalties reduced the priority of global yield or biomass outputs, canopy-only remote sensing, non-fruit targets, and image-level outputs. The score did not decide inclusion.

To prevent one model family from dominating, selection used buckets for core identity or inventory, direct oil palm evidence, fruit multiview or 3D evidence, fruit-instance baselines, transfer mechanisms, and prior reviews or positioning. The resulting 250-record shortlist and 60-record first wave are machine-readable in \path{literature/search/derived/priority-shortlist-2026-08-10.csv} and \path{literature/search/derived/priority-review-wave1-2026-08-10.csv}. This design makes the next review step reproducible without pretending that ambiguous candidates have already been included or excluded.

\subsection{Two evidence registers}

The manuscript keeps two registers separate. Register A is the new Scopus and OpenAlex search, including all candidate and screening states. Register B is the targeted evidence register: studies whose full text was retrieved and whose dataset, protocol, quantitative evaluation, and mechanism fields were verified in the ledger. The ranking and shortlist are selection aids inside Register A; they are not inclusion decisions. This separation prevents the 20,035 search pool from being presented as 20,035 included studies.

\section{The inventory task as a design space}
\label{sec:designspace2}

\subsection{From detections to inventory}

The task has five distinct operations. First, a detector or segmenter proposes an instance in each observation. Second, a class head attaches an attribute to that instance. Third, observations are associated when they refer to the same physical object. Fourth, duplicate observations are rejected or merged. Fifth, the remaining identities are aggregated into counts per class. The first two operations can be evaluated with precision, recall, average precision, and class confusion. The last three require identity and count metrics.

This decomposition changes the interpretation of a strong detector. A detector with high average precision on independent images may still overcount when a tree is filmed from both sides. A tracker may preserve identity through a short video but fail when the camera returns after a gap. A 3D association method may remove duplicates but become unusable when depth is uncalibrated. The design choice is therefore conditional on the acquisition regime.

\subsection{Six identity mechanisms}

Table~\ref{tab:mechanisms2} defines the six mechanisms used in the synthesis. M0 is the no-association baseline. M1 uses a statistical correction or a fixed counting protocol without assigning a persistent identity. M2 matches local appearance or a learned re-identification embedding. M3 uses temporal continuity, motion, and track management. M4 uses calibrated geometry, depth, SfM, or a 3D map. M5 learns a multi-view association function directly from linked observations. A system can combine mechanisms, but its reported result must identify which mechanism is responsible for duplicate rejection.

\begin{table}[t]
\centering
\caption{Identity mechanisms and their minimum assumptions.}
\label{tab:mechanisms2}
\small
\begin{tabularx}{\columnwidth}{@{}lY Y@{}}
\toprule
Code & Mechanism & Main assumption and measurable risk\\
\midrule
M0 & Intra-view detection or segmentation & One view is sufficient. Hidden objects and repeated views are not resolved.\\
M1 & Statistical correction or fixed census & Duplicate structure is stable enough for a correction factor. The correction may not transfer across trees.\\
M2 & Appearance or re-identification & The instance retains discriminative texture, colour, or shape despite illumination and occlusion.\\
M3 & Temporal tracking & Consecutive observations overlap in time and motion is sufficiently smooth. Track fragmentation and drift remain possible.\\
M4 & Geometric or 3D association & Calibration, scale, or reliable reconstruction is available. Sensor failure and SfM cost are risks.\\
M5 & Learned multi-view association & Linked training data represent view, class, and occlusion variation. The model may overfit one acquisition pattern.\\
\bottomrule
\end{tabularx}
\end{table}

\subsection{Acquisition assumptions}

The second design axis is the assumption set. A1 is the amount of overlap and occlusion. A2 is the camera baseline or view separation. A3 is motion smoothness or scene stationarity. A4 is calibration and metric scale. A5 is appearance distinctiveness. A6 is regularity of duplicate exposure, for example a fixed four-side census. A7 is the availability of identity supervision or an oracle association. M1 mainly depends on A6. M2 depends on A5. M3 depends on A3. M4 depends on A2 and A4. M5 depends on A7 and on coverage of the training acquisition pattern.

This two-axis view gives a precise meaning to the phrase ``multi-view deduplication.'' It is not a mandatory module. It is a mechanism whose value should be tested against the assumptions that survive in the field. A fixed four-side protocol may make M1 competitive when duplication is regular and metric depth is unreliable. An unstructured moving camera may require M3 or M4 because the number and order of observations are not fixed.

\subsection{Class attributes beyond maturity}

Maturity classification is common in FFB literature, but the identity dependency is more general. An instance may carry a maturity label, a size class, a defect label, a disease state, a cultivar label, or a harvest-readiness state. The inventory should first establish one identity record and then attach one or more attributes to that record. If the same bunch appears in three views with two maturity predictions, the system must report whether the conflict is resolved by confidence, temporal evidence, geometry, or manual adjudication. Summing per-view class predictions without this step is not a class-wise inventory.

\section{Direct agricultural and fruit evidence}
\label{sec:direct2}

\subsection{What the oil-palm literature measures}

The direct oil-palm literature is concentrated on detection, ripeness classification, dataset construction, and yield-related counting. Lai et al. reviewed FFB ripeness detection methods across sensing families, while Goh et al. reviewed non-destructive classification methods including computer vision, LiDAR, spectroscopy, thermal sensing, and laser-light backscatter \cite{lai2023palmslr,goh2025palmreview}. These reviews establish that sensor choice and class definition are central problems. Their primary unit is ripeness assessment, not the persistent identity of one bunch across repeated views.

Naftali et al. compared YOLO and other detectors on an oil-palm plantation dataset and reported a depthwise-separable YOLOv8 variant as a real-time candidate \cite{naftali2024palmcounter}. This is useful evidence for M0 detector selection and deployment constraints. It does not by itself measure whether a bunch detected in one frame is the same bunch detected in another frame. A detector benchmark and an inventory benchmark answer different questions.

Three dataset papers expose the acquisition variables needed for the new design space. Goh et al. released 400 RGB images and 400 registered depth maps from four plantation locations, captured with an Intel RealSense D455f. The article reports a mean registration error of 1.8 cm at 3 m and 92.5\% expert inter-rater agreement for binary ripeness labels \cite{goh2025palmdataset}. Suharjito et al. released plantation imagery captured from multiple angles with five categories: unripe, underripe, ripe, flower, and abnormal; the dataset contains 10,207 training images, 2,896 validation images, and 1,400 test images after the reported augmentation and split procedure \cite{suharjito2025palmdataset}. These datasets address sensor and class variation, but their public descriptions do not make the same cross-view identity graph the primary evaluation object.

SawitMVC is directly aligned with the inventory problem. Its source article describes 3,992 smartphone images of 953 oil-palm trees from two plantations, captured from four or eight angular positions. The dataset uses four BBC maturity classes and a separate per-tree JSON layer with 9,823 unique bunches and cross-view identity links \cite{indriani2026sawitmvc}. This makes a tree-level unique-count metric possible. It also exposes the central confounder: the number of visible bounding boxes is not the number of unique bunches.

Semantic NeRF converts unordered posed RGB observations into a 3D fruit field and clustered fruit point cloud, providing a geometric route to consolidate repeated observations before counting \cite{ledger_m_1210e3ea00486b68}.

The tomato RGB-D SLAM study projects fruit masks into a filtered semantic point cloud and fits spheres to estimate fruit counts and volume at plant and row levels \cite{ledger_r_b2e0e9aabad17a97}.

A pear study fuses RGB detections, LiDAR, FAST-LIO2 SLAM, and temporal association for 3D fruit counting; field validation reports 96.2 percent counting accuracy while retaining 53 double-counts linked to association threshold and SLAM drift \cite{ledger_r_d4270ea8ab11f268}.

\subsection{Fruit and orchard evidence}

Single-view fruit-load studies show the limit of a detector-only interpretation. Koirala et al. compared six existing architectures and developed MangoYOLO for mango detection. Their best reported model achieved an F1 score of 0.968 and average precision of 0.983 on an independent test set, while orchard load estimates still required a correction factor against hand-harvest counts \cite{koirala2019deepfruit}. The correction is an M1 mechanism: it can be operationally useful without solving identity, but its validity depends on stable exposure and sampling assumptions.

Gené-Molá et al. made identity and 3D location explicit by combining Mask R-CNN instance segmentation with SfM, projecting image detections into a 3D point cloud, and removing false positives with a support vector machine. The study used 11 Fuji apple trees containing 1,455 apples. The reported F1 score increased from 0.816 for 2D detection to 0.881 for 3D detection and location, while the authors identified SfM processing time as a barrier to real-time use \cite{genemola2020fruit3d}. This is direct evidence that geometric association can change the counting result, but it is not evidence that SfM is the best field solution for oil palm.

Fu et al. used RGB and depth features with Faster R-CNN for apple detection in dense foliage and robotic harvesting \cite{fu2020apple}. The study supports the role of depth as an occlusion and localization cue. It does not settle whether depth remains reliable under bright outdoor oil-palm conditions. The new search also returned citrus and grape yield-estimation studies, including multi-period citrus estimation and a YOLO-based grape framework \cite{zhai2025citrusyields,ma2026grapeyield}. These records widen the agricultural comparison, but their exact full-text inclusion status remains tied to the retrieval boundary in Section~\ref{sec:protocol2}.

\subsection{Positioning against previous reviews}

The reviews found by Q5 are adjacent rather than identical to this review. Table~\ref{tab:position2} records the scope visible in the accessible article record and identifies whether cross-view identity is treated as the primary unit. The wording is deliberately bounded: ``not explicit'' means that the accessible scope and abstract do not make identity and duplicate rejection the central organizing question. It does not claim that no cited paper inside the review discusses tracking.

\begin{table}[t]
\centering
\caption{Positioning against previous review and survey records.}
\label{tab:position2}
\scriptsize
\begin{tabularx}{\columnwidth}{@{}p{0.25\columnwidth}p{0.21\columnwidth}p{0.18\columnwidth}Y@{}}
\toprule
Review & Primary scope & Cross-view identity central? & Position relative to this review\\
\midrule
Lai et al. \cite{lai2023palmslr} & Oil-palm FFB ripeness detection methods & Not explicit & Supplies direct agricultural sensing context; does not organize unique counting mechanisms.\\
Goh et al. \cite{goh2025palmreview} & FFB ripeness classification across non-destructive sensors & Not explicit & Supplies modality and deployment limitations; this review adds identity and inventory evaluation.\\
Naftali et al. \cite{naftali2024palmcounter} & Oil-palm detector and counter comparison & Not explicit & Supplies detector and edge-computation context; this review separates detector score from duplicate-free count.\\
Naghipour et al. \cite{naghipour2024fruitreview} & AI methods for fruit detection and classification & Not explicit & Supplies broad fruit and model coverage; this review makes cross-view association a design axis.\\
Sapkota et al. \cite{sapkota2024yoloagri} & YOLO variants in agricultural applications & Not explicit & Supplies architecture lineage and application taxonomy; this review is task-first rather than model-first.\\
Nordin and Misro \cite{nordin2026palmslr} & Deep-learning review of palm-oil maturity detection & Not explicit in accessible record & Adds a recent maturity-focused comparator; this review generalizes the class attribute beyond maturity.\\
\bottomrule
\end{tabularx}
\end{table}

The bounded conclusion from Q5 is therefore not ``no previous review exists.'' It is that the accessible review set is organized mainly by sensor, ripeness, detector family, or agricultural application. The cross-view identity of each physical fruit as the prerequisite for a class-wise inventory is not the primary comparison unit in that set. This is the gap addressed here, subject to the stated search and retrieval boundary.

\section{Transferable mechanisms for identity and duplicate rejection}
\label{sec:transfer2}

\subsection{Appearance and temporal association}

M2 matches instances by appearance. A descriptor may use colour, texture, shape, or a learned embedding. Its strength is that it does not require metric depth. Its failure modes are predictable in fruit scenes: two bunches of the same maturity can look alike, illumination changes colour, and partial occlusion removes the most distinctive region. An evaluation must therefore report identity precision and recall under appearance changes, not only detector average precision.

M3 uses time. A tracker predicts where an instance should appear in the next frame, associates detections, and terminates a track when evidence is insufficient. Video-based oil-palm ripeness work shows why temporal input is attractive for deployment \cite{junior2023videopalm}. Temporal continuity is weaker when the camera moves around a tree and revisits a previously observed side. The report must state the maximum time gap and whether the camera path is ordered or arbitrary.

\subsection{Geometry and 3D association}

M4 associates observations through camera pose, depth, point clouds, or an SfM reconstruction. Gené-Molá et al. provide the clearest fruit example: a 2D instance mask is projected into a shared 3D representation so that repeated image detections can be consolidated \cite{genemola2020fruit3d}. Park et al. describe a newer real-time orchard approach that incrementally builds 3D crop instances and uses 3D generalized intersection-over-union for association \cite{park2026cropassociation}. The mechanism is directly aligned with duplicate rejection, but its field cost includes calibration, reconstruction stability, and sensor coverage.

Non-agricultural 3D detection studies provide transferable design components rather than fruit results. Camera and LiDAR fusion methods such as 3D-CVF and TransFusion show how multiple sensor views can be aligned at feature or decision level \cite{yoo20203dcvf,bai2022transfusion}. Visual SLAM systems show how a moving camera can estimate pose and maintain a scene map, while dynamic-scene variants explicitly mask moving objects before mapping \cite{bescos2018dynaslam}. These studies motivate a measurement question for agriculture: does a shared geometric frame reduce false duplicates after controlling for detector recall and sensor failure?

\subsection{Learned multi-view association}

M5 learns identity from linked observations. Training examples must state which detections belong to one physical object and which are different objects. Without this supervision, a model can learn a view-specific shortcut and still appear accurate on a random image split. SawitMVC's per-tree identity links create the type of supervision required for this experiment \cite{indriani2026sawitmvc}. The necessary split is by physical tree or plantation block, not by randomly distributing frames from one tree across train and test.

\section{Modalities and fusion choices}
\label{sec:fusion2}

RGB provides colour, texture, and fine appearance. Its weaknesses are shadows, specular highlights, camouflage, and scale changes. Depth provides a spatial cue, but ``depth'' is not one measurement. A calibrated sensor may provide metric range; stereo and SfM provide geometry relative to a reconstructed scene; a monocular network provides a learned or pseudo-depth ordering whose scale may not be metric. A manuscript must state which meaning is used.

Three fusion points are useful design categories. Early fusion concatenates RGB and depth at the input. It is simple but exposes the detector to missing or misregistered channels. Middle fusion extracts separate features and combines them at an intermediate layer, which permits modality-specific processing and quality gating. Late fusion combines detections, scores, or identity graphs after each modality has produced an estimate. RGB-depth object-detection studies demonstrate that the choice of fusion point is a measurable hypothesis rather than a universal rule \cite{ophoff2019multimodal}. The oil-palm RGB-depth CNN and visual-SLAM system illustrates a direct agricultural attempt to combine appearance, depth, and localization \cite{siong2021rgbdepthpalm}.

The design-space consequence is that sensor quality belongs in the evaluation, not only in the apparatus description. Each experiment should report the valid-depth ratio, RGB-depth registration error, missing-depth policy, and the fraction of detections supported by valid geometry. A fusion gain observed after dropping invalid frames is not comparable with a gain measured on all frames. The Goh dataset's reported registration error provides an example of the type of calibration evidence needed before interpreting a depth-based identity result \cite{goh2025palmdataset}.

\section{Synthesis: from detector score to inventory evidence}
\label{sec:synthesis2}

The evidence supports five design rules.

\begin{enumerate}[leftmargin=*,itemsep=2pt]
  \item Report per-view detection separately from cross-view identity. A detector score cannot establish duplicate-free counting.
  \item Define the unit of truth at tree level. For each physical tree, store the unique instance count, per-view appearances, class labels, and visibility state.
  \item Treat the acquisition path as an experimental factor. A fixed four-side census, a continuous video, and an arbitrary multi-camera setup impose different assumptions.
  \item Test identity mechanisms against their required assumptions. M1 needs regular duplication; M3 needs temporal overlap; M4 needs geometry; M5 needs identity-linked training data.
  \item Report class-wise count error after identity resolution. A correct class label on a duplicate is still a wrong inventory.
\end{enumerate}

The minimum evaluation bundle is therefore: per-instance precision and recall, class-wise average precision, identity precision and recall, duplicate rate, missed-appearance rate, absolute and relative count error per tree, class-confusion error per unique instance, sensor-validity statistics, and latency. Where tracks are used, IDF1 or HOTA can complement track precision and recall. Where 3D association is used, calibration error and association distance must be reported. These metrics should be paired by physical tree so that a difficult canopy is not assigned to one method and an easy canopy to another.

\subsection{Focused evidence matrix}

Appendix~\ref{app:matrix3} contains one row per the 44 studies included after targeted full-text review. The table records method, modality, evaluation context, evidence pages, identity mechanism, class attribute, unique-instance ground truth status, duplicate handling, finding, and limitations. The machine-readable matrix is \path{audit/evidence-matrix-v2.csv}; the assumptions and violated-condition codes are explicitly marked as reviewer inferences.

\section{Research gaps and testable agenda}
\label{sec:agenda2}

The first gap is an evaluation gap. Many studies report detection or ripeness classification, but a class-wise inventory requires identity links and a tree-level unique-count ground truth. SawitMVC demonstrates the data structure needed for this measurement, while the earlier fruit and FFB literature often reports image-level or frame-level results \cite{indriani2026sawitmvc,naftali2024palmcounter,koirala2019deepfruit}.

The second gap is a mechanism comparison gap. The literature contains examples of statistical correction, temporal video processing, RGB-depth fusion, and SfM, but these mechanisms are rarely compared under one acquisition protocol with the same detector and the same tree-level truth. The appropriate experiment is not ``model A versus model B'' alone. It is an M0-to-M4 or M5 ablation in which the detector, image split, class labels, and counting truth remain fixed.

The third gap is a sensor-assumption gap. A calibrated depth camera, a point cloud, LiDAR, and monocular pseudo-depth should not be grouped under one label. Each has a different scale, failure mode, and deployment cost. An outdoor FFB benchmark should report invalid depth, registration, range, illumination, and view baseline before assigning a performance gain to the modality.

The fourth gap is a class-resolution gap. Maturity is common, but the same identity record may need multiple attributes. A future benchmark should permit at least one class attribute beyond maturity or make clear why a single attribute is the intended use case. The evaluation should distinguish class confusion from identity duplication.

The resulting research agenda is a sequence of gates. Gate 1 verifies the per-view detector and the class labels. Gate 2 builds the unique-instance tree truth and measures the M0 baseline. Gate 3 tests M1 under the fixed view protocol. Gate 4 adds temporal or appearance association. Gate 5 adds calibrated geometry or registered depth. Gate 6 tests a learned association model with a tree-disjoint split. A mechanism advances only if it lowers tree-level count error without exceeding the latency and sensor-quality budget defined before the experiment.

\section{Limitations}
\label{sec:limitations2}

The search is reproducible from the captured raw exports and scripts, but the full-text stage is intentionally targeted rather than exhaustive. Scopus was available and Web of Science was not, so database coverage is not equivalent to a two-subscription search. OpenAlex metadata and abstracts are useful for discovery but do not replace article full text. The local audit found six matches in the legacy PDF corpus; targeted retrieval then produced 44 complete evidence records. Consequently, the 20,035 records remain a ranked search pool, not an included-study count.

The focused matrix also mixes direct agricultural studies, fruit studies, transferable mechanism papers, and reviews. This is deliberate for a design-space review, but the evidence classes are not interchangeable. A result from an autonomous-driving benchmark does not establish oil-palm performance. A review that discusses ripeness does not establish cross-view identity. A dataset descriptor establishes the availability of labels and sensors, not the superiority of a model.

The review does not estimate a pooled accuracy, rank all detectors, or claim that explicit deduplication is always beneficial. Its bounded claim is that cross-view identity is a necessary design variable for the stated inventory task and that the value of each identity mechanism should be evaluated under its acquisition assumptions.

\section{Conclusion}
\label{sec:conclusion2}

The target of agricultural multi-view perception is a unique, class-wise inventory, not a sum of independent detections. The revised protocol makes the search reproducible and makes the selection boundary explicit: 20,035 candidates were ranked, 250 were shortlisted, and 44 targeted full-text studies supplied the current evidence matrix. The direct evidence includes oil-palm detector and dataset studies, multi-angle and RGB-depth systems, fruit SfM, NeRF, LiDAR-camera fusion, and orchard counting. Together they support a six-mechanism design space spanning intra-view counting, statistical correction, appearance, temporal tracking, geometry, and learned association.

The central contribution is a testable question rather than a fixed architecture: when does explicit identity resolution reduce tree-level duplicate and class-wise count error enough to justify its assumptions and cost? Answering that question requires tree-disjoint data, identity-linked annotations, sensor-quality reporting, and evaluation after aggregation. Oil palm is the principal case, but the design logic applies to other tree-fruit inventories whenever repeated observations must be converted into one count.

\appendix

\section{Focused evidence matrix}
\label{app:matrix3}

The matrix contains one row per study included after targeted full-text review. Direct evidence covers agriculture and fruit systems; transferable evidence answers an identity or geometry mechanism question outside agriculture. The assumptions and violated-condition codes are reviewer inferences, not verbatim claims from every source. The machine-readable version is available at \path{audit/evidence-matrix-v2.csv}.

\small
\setlength{\LTleft}{0pt}
\setlength{\LTright}{0pt}
\setlength{\tabcolsep}{2pt}
\makeatletter
\@ifclassloaded{IEEEtran}{\onecolumn}{ }
\makeatother
\begin{longtable}{@{}>{\raggedright\arraybackslash}p{0.14\textwidth}>{\raggedright\arraybackslash}p{0.18\textwidth}>{\raggedright\arraybackslash}p{0.18\textwidth}>{\raggedright\arraybackslash}p{0.16\textwidth}>{\raggedright\arraybackslash}p{0.25\textwidth}>{\raggedright\arraybackslash}p{0.07\textwidth}@{}}
\caption{Targeted row-level evidence matrix. Pages refer to the verified full text recorded in the ledger.}\label{tab:matrix3}\\
\toprule
Study & Method and modality & Evaluation and pages & Identity and attribute & Finding and limitation & Status\\
\midrule
\endfirsthead
\multicolumn{6}{c}{\tablename\ \thetable{} continued}\\
\toprule
Study & Method and modality & Evaluation and pages & Identity and attribute & Finding and limitation & Status\\
\midrule
\endhead
\midrule
\multicolumn{6}{r}{Continued on next page}\\
\endfoot
\bottomrule
\endlastfoot
Von Hirschhausen, Laura-Sophia (2025) \cite{ledger_m_3f8f02033fe9db3a} & calibrated\_\allowbreak{}stereo\_\allowbreak{}rgb\_\allowbreak{}multiview; apple & object\_\allowbreak{}detection\_\allowbreak{}F1\_\allowbreak{}AP\_\allowbreak{}AR\_\allowbreak{}mAP, BBCH\_\allowbreak{}accuracy, qualitative\_\allowbreak{}3D\_\allowbreak{}reconstruction Pages 1,2,3,4,5,6,7. & M4; none; GT n; dedup no & The dataset supports fruit detection, phenological stage classification, and calibrated stereo reconstruction with quantitative detection metrics. Limit: Per-instance fruit boxes and stereo views are explicit, but cross-view unique identity or deduplication is not evaluated, do not treat this paper as proof of duplicate-free inventory. & direct\\
Fusaro, Daniel (2026) \cite{ledger_r_0eda6374aaf85fae} & colored\_\allowbreak{}3d\_\allowbreak{}point\_\allowbreak{}cloud\_\allowbreak{}lidar; agricultural fruit or transferable mechanism & PQ\_\allowbreak{}SQ\_\allowbreak{}RQ\_\allowbreak{}IoU, F1p\_\allowbreak{}F1n\_\allowbreak{}mF1, five\_\allowbreak{}fold\_\allowbreak{}cross\_\allowbreak{}validation, runtime Pages 1,2,5,6,7,9,12,13,14. & M2/M3/M4; none; GT n; dedup explicit & Learned fruit descriptors and attention-based matching improve temporal identity association over the compared baselines on the strawberry benchmark, while the instance segmentation module also performs across strawberry and apple datasets. Limit: Re-identification is evaluated only on strawberries because temporally consistent instance annotations are unavailable for other fruits, and segmentation errors propagate into re-identification. The sensor is a high-resolution terrestrial LiDAR scanner. & direct\\
Suchet Bargoti (2016) \cite{ledger_r_143d70deec8d2b1d} & monocular\_\allowbreak{}rgb\_\allowbreak{}ugv; apple & pixel\_\allowbreak{}F1, detection\_\allowbreak{}F1, count\_\allowbreak{}r\_\allowbreak{}squared, yield\_\allowbreak{}r\_\allowbreak{}squared, yield\_\allowbreak{}error Pages 1,2,3,4,12,16,17,18,19,24. & M0; none; GT n; dedup no & CNN segmentation followed by Watershed detection produced the reported best detection F1 of 0.858 and row-yield r-squared of 0.826 on the apple orchard data. Limit: The study counts detected fruits within images and aggregates rows, but it does not evaluate persistent identity or duplicate rejection across repeated views of one tree. & direct\\
Stein, Madeleine (2016) \cite{ledger_r_408d13511cac4177} & rgb\_\allowbreak{}gps\_\allowbreak{}ins\_\allowbreak{}lidar; mango & count\_\allowbreak{}r\_\allowbreak{}squared, slope, percentage\_\allowbreak{}error, repeatability, triangulation\_\allowbreak{}nearest\_\allowbreak{}neighbour\_\allowbreak{}distance Pages HTML abstract, sections 1, 2, 3.3, 3.4, 3.5, 4.1, 4.4, 5. & M3/M4; none; GT n; dedup implicit & Multi-view tracking with epipolar geometry and tree masks converts repeated mango observations into tree-linked counts with near-unity calibration slope and 1.36 percent over-count on the 16-tree validation set. Limit: The evaluation is on mango orchard traversal data and depends on navigation quality, timing, image coverage, and LiDAR masks. The authors report that transfer to other fruit varieties or pruning regimes requires further testing. & direct\\
Liu, Xu (2018) \cite{ledger_m_9966d81433d61ed3} & monocular\_\allowbreak{}rgb\_\allowbreak{}video\_\allowbreak{}sfm; agricultural fruit or transferable mechanism & count\_\allowbreak{}ground\_\allowbreak{}truth, L1\_\allowbreak{}loss, error\_\allowbreak{}mean, error\_\allowbreak{}standard\_\allowbreak{}deviation, before\_\allowbreak{}after\_\allowbreak{}correction Pages 1,2,3,4,5,6,7,8. & M3/M4; none; GT n; dedup explicit & The 3D reconstruction correction reduces duplicate and outlier tracks and substantially improves count error on both orange and apple sequences. Limit: The method uses relative 3D localization from monocular sequences and heuristic size and depth thresholds. Performance depends on illumination, occlusion, camera stability, and track continuity. & direct\\
Fazida Hanim Hashim (2018) \cite{ledger_r_97003608b250cdc5} & lidar\_\allowbreak{}intensity\_\allowbreak{}point\_\allowbreak{}cloud\_\allowbreak{}905nm; oil palm FFB & intensity\_\allowbreak{}ranges, reflectance\_\allowbreak{}bands, four\_\allowbreak{}tree\_\allowbreak{}field\_\allowbreak{}scan, fruit\_\allowbreak{}to\_\allowbreak{}leaf\_\allowbreak{}ratio Pages 1,4,5,6,8. & M0; maturity; GT n; dedup no & LiDAR intensity point clouds provide a preliminary non-destructive way to separate three FFB ripeness levels and support field-level fruit-to-leaf observations. Limit: The paper does not report a conventional held-out classification accuracy or cross-view identity metric. The experiment uses a small number of fruitlets per category and four field-scanned trees. & direct\\
Fawcett, Dominic (2019) \cite{ledger_m_084cfd5badbb63d3} & UAV\_\allowbreak{}RGB\_\allowbreak{}SfM\_\allowbreak{}point\_\allowbreak{}cloud; oil palm FFB & 98.2\% mapping accuracy for 776 palms; stem-height RMSE 0.27 m (R2=0.63) for seven-year palms and 0.45 m (R2=0.69) for ten-year palms; mean point-cloud precision 26.7 cm Pages PDF pp. 7538-7560; abstract; sections 2.1-2.2.2, 2.4, 3.1, 5. & M4; size; GT y; dedup explicit & UAV SfM canopy-height models can support per-palm inventories and structural measurements, providing a transfer mechanism for geometry-based object association while not proving fruit-level cross-view deduplication. Limit: Tree-canopy study rather than fruit or FFB inventory; local-maxima errors increase with undergrowth and overlapping canopies; no explicit cross-view fruit identity metric or duplicate-rejection experiment. & direct\\
Liu, Xu (2019) \cite{ledger_r_00f82e0f404d691b} & monocular\_\allowbreak{}RGB\_\allowbreak{}video; mango & 18 ground-truth trees; average 175 mangoes per tree; Mono Final mean error 27.8 and standard deviation 29.6 fruits; regression slope 0.86 and R2=0.78, compared with sensor-suite multi-view slope 0.97 and R2=0.88 Pages pp. 1, 3-8; sections IV-VI. & M3/M4; maturity; GT y; dedup explicit & A monocular RGB sequence can support duplicate-aware fruit inventory by combining tracking, semantic SfM, landmark re-association, and tree-centroid depth filtering across opposite views. Limit: Undercounting increases with heavy occlusion and high fruit counts; evaluation uses 18 manually counted trees and manual target-tree masks; performance depends on sequence ordering and does not classify fruit maturity or other class attributes. & direct\\
Häni, Nicolai (2020) \cite{ledger_r_1358c0868a62c84a} & RGB\_\allowbreak{}image\_\allowbreak{}sequences\_\allowbreak{}plus\_\allowbreak{}semantic\_\allowbreak{}mapping; apple & Seven detection test datasets and three two-side yield datasets; merged CNN counting yield 95.56\%-97.83\% of harvested ground truth; GMM detection highest F1 on six of seven datasets; CNN cluster counting reported about 90.5\% accuracy Pages pp. 1-6, 11-16, 20-25; sections 3.3, 4, 5, 6.4, 7. & M3/M4; maturity; GT n; dedup implicit & This is a strong design-space benchmark showing that fruit inventory quality depends on the composition of per-frame detection, cluster counting, temporal tracking, and geometric merging of opposite row sides, rather than detector choice alone. Limit: Evaluation is concentrated on apple orchard row datasets; occlusion, changing illumination, annotation noise, and side-specific reconstruction affect performance; the row-merging mechanism is geometric and does not provide general fruit re-identification across arbitrary scenes or class attributes beyond counting. & direct\\
Alfatni, Meftah Salem M. (2020) \cite{ledger_r_2f7d604b1d3fa5e4} & controlled\_\allowbreak{}RGB\_\allowbreak{}CCD; oil palm FFB & Three ROI sizes and colour feature families evaluated with ANN, KNN, and SVM; best reported ANN histogram result 94\% test accuracy and 1.6 s processing at ROI3 100x100; AUC values above 90\% Pages article pages 1-9; sections 2, 3.1, 3.2, 4; Tables 1-2. & M0; maturity; GT n; dedup no & Provides direct oil-palm class-attribute evidence for maturity grading and shows how an instance-level FFB attribute can be attached after segmentation, but it does not address duplicate-free inventory across tree observations. Limit: Controlled conveyor and illumination setting; single-bunch image grading rather than orchard or tree-level multi-view inventory; no persistent identity, cross-view association, or duplicate-rejection evaluation. & direct\\
Husin, Husna Sarirah (2021) \cite{ledger_r_3e45c6dc2566f46b} & LiDAR\_\allowbreak{}intensity\_\allowbreak{}point\_\allowbreak{}cloud\_\allowbreak{}and\_\allowbreak{}latitude\_\allowbreak{}longitude; oil palm FFB & Three overripe, six ripe, and five under-ripe dataset plots are discussed; mean-intensity ranges and threshold table are reported, but no classification accuracy or held-out test metric is provided Pages article pages 1-9; sections III-V; Table I. & M4; maturity; GT n; dedup no & The study links a class attribute to a spatially localized FFB record, showing a possible inventory record structure, but spatial mapping alone is not evidence that repeated observations of the same FFB have been deduplicated. Limit: Variable LiDAR resolution can return intensity from non-FFB surfaces; sample collection was insufficient for a high-accuracy system and was not performed at the target plantation field because of COVID-19; no cross-view association, persistent identity, or accuracy metric. & direct\\
Lai, Jin Wern (2022) \cite{lai2022yolov4palm} & RGB\_\allowbreak{}depth\_\allowbreak{}camera\_\allowbreak{}RealSense\_\allowbreak{}D435; oil palm FFB & 49 sampled trees; 240 positive and 250 negative images; 20-tree field test; mAP 87.9\%, recall 82\%, F1 88\%, mean IoU 70.19\%, and approximately 21 FPS Pages author version pp. 1-5; sections II-IV; Tables 1-6. & M0; maturity; GT n; dedup no & Per-instance FFB detection can attach a ripeness class to a live tree observation and trigger a harvesting action, but the study does not establish whether the same FFB seen in multiple frames is counted once. Limit: Only ripe versus unripe detection; no overripe class, fruit count, persistent ID, cross-view association, or duplicate-rejection metric; field test used 20 similar-age trees under a sunny morning condition and manual labels. & direct\\
Mansour, Mohamed Yasser Mohamed Ahmed (2022) \cite{ledger_r_19d96439ecdda18d} & RGB\_\allowbreak{}ground\_\allowbreak{}FFB\_\allowbreak{}images; oil palm FFB & 304 balanced images; 70/20/10 train-validation-test split; YOLOv5m mAPval 0.842 at IoU 0.5:0.95; MobileNetV2 SSD 0.478; EfficientDet-lite0/1/2 0.743/0.803/0.812; YOLOv5n/s 0.781/0.832 Pages pp. 1326-1335; sections 3.1, 4, 5; Table 2. & M0; maturity; GT n; dedup no & This is direct oil-palm evidence that a per-instance detector can attach a broader maturity class set than binary ripe versus unripe, but it remains an intra-view classification result rather than duplicate-free inventory. Limit: All images were ground FFBs after harvest due COVID-19 access restrictions, not on-tree observations; dataset is small and balanced at 304 images; no multi-view identity, persistent fruit ID, counting, or duplicate-rejection evaluation; authors identify on-tree and diverse-field imagery as future work. & direct\\
Meili Sun (2022) \cite{ledger_r_1a70fa2aeed54ea3} & RGB\_\allowbreak{}nighttime\_\allowbreak{}images\_\allowbreak{}LED\_\allowbreak{}and\_\allowbreak{}smartphone\_\allowbreak{}Canon; apple & NightFruit: 1960 augmented images, 1371 train and 589 test, 7713 and 3396 fruits, AP 85.2 percent, APs 67.5 percent, AP50 97.8 percent; Gala: 1361 images, 953 train and 408 test, 4943 and 2194 fruits, AP 61.0 percent, APs 45.2 percent; corruption robustness reported across fog, elastic, brightness, noise, frost, and motion blur. Pages pp.4421-4432; sections 2,4.2,4.5,4.7,5,6. & M0; maturity; GT y; dedup no & Nighttime small-green-fruit detection can provide useful per-image fruit observations under occlusion and variable weather, but detector outputs alone do not solve temporal tracking, cross-view identity, or unique inventory counting. Limit: Random image splits do not prove scene or tree-disjoint generalization; no temporal tracking, re-identification, or cross-view deduplication; performance degrades under brightness, motion blur, and heavy occlusion; Gala AP remains modest. & direct\\
Lu, Jianqiang (2022) \cite{ledger_r_2828a3937318ca86} & RGB\_\allowbreak{}handheld\_\allowbreak{}and\_\allowbreak{}cloud\_\allowbreak{}camera; citrus & 557 original images, 3273 annotated fruit instances, visible targets with more than 70 percent occlusion ignored, reported 6:2:2 train validation test split after augmentation; improved Mask-RCNN plus ResNet50 AP50 95.36 percent, AP75 93.45 percent, PR area 0.9673, field inference about 0.939 to 0.969 seconds on AGX Xavier. The article reports inconsistent augmented/test image counts in different sections, which is retained as a limitation. Pages pp.1-5,10-16; sections 2.1-2.6, 3, 4, 5, 6; Tables 1-3. & M0; none; GT y; dedup no & CB-Net feature fusion improves per-image citrus green-fruit instance segmentation in complex green backgrounds and supports orchard thinning decisions, but the study remains an intra-view detector and does not produce unique cross-view fruit inventories. Limit: No temporal tracking, re-identification, or cross-view deduplication; random image split does not establish tree-disjoint or orchard-disjoint generalization; heavily occluded fruit is ignored during annotation; overlap remains a stated failure mode; only sugar tangerine and one green-fruit target class are evaluated; reported augmented dataset counts are internally inconsistent. & direct\\
Wan Nurazwin Syazwani, R. (2022) \cite{ledger_r_3f230a23a2a956c8} & UAV\_\allowbreak{}top\_\allowbreak{}view\_\allowbreak{}RGB; apple & DJI Phantom 3 Advanced, 1300 frames from 110 rows over 1.44 acres, 20 manually annotated sample images, 180 pineapple and 180 background instances, 90 fruit and 90 noise used for train and test, six classifier families and ANN training variants; best ANN-GDX with ANOVA reported 94.4 percent accuracy and 87 of 90 pineapple crowns correctly counted. Pages pp.1265-1276; sections 2.1-2.8, 3, 4; Figs.1-9. & M0; maturity; GT n; dedup no & Top-view UAV RGB processing can identify and count pineapple crowns using engineered colour, shape, and texture features, providing a transferable agricultural example of per-object detection followed by plot-level counting. Limit: No temporal tracking, re-identification, or cross-view deduplication; the evaluation uses only 20 sampled images and 180 fruit instances from one plantation division; train-test separation at instance level does not establish scene or row-disjoint generalization; crown detection is a proxy for fruit and no fruit maturity attribute is modeled. & direct\\
Yang, Lu (2022) \cite{ledger_r_4b3d8d423dbd5d0e} & RGB\_\allowbreak{}multi\_\allowbreak{}camera\_\allowbreak{}vehicle\_\allowbreak{}images; non-agricultural transfer domain & VeRi-776, VehicleID, VRIC, and VRAI benchmarks with CMC and mAP; baseline plus WCVL gains include VehicleID small mAP 86.8 to 90.3 percent, medium 83.6 to 87.1 percent, large 81.0 to 84.6 percent, VRIC CMC1 74.3 to 76.2 percent, and VRAI mAP 74.4 to 80.1 percent. Pages pp.1-9; sections I-VI; Tables I-VIII. & M2; none; GT y; dedup n-a & Weakly supervised cross-view embeddings can reduce viewpoint-induced appearance variation using identity labels only, offering a transferable mechanism for associating repeated observations when viewpoint labels are unavailable. Limit: The benchmark concerns vehicles rather than agricultural objects; it assumes identity labels during training; failures remain under extreme brightness and near-identical appearance; it does not address fruit occlusion, tree structure, changing object appearance, or agriculture-specific instance inventory. & transferable\\
Zhang, Chenglong (2022) \cite{ledger_r_c79d61a639a09929} & UAV\_\allowbreak{}and\_\allowbreak{}ground\_\allowbreak{}vehicle\_\allowbreak{}RGB\_\allowbreak{}SfM\_\allowbreak{}point\_\allowbreak{}clouds\_\allowbreak{}GPS; apple & Dutch Elstar orchard with 14 rows and 100 trees in the monitored row; manual ground truth for 31 random trees; segmentation validation on 107 trees and 186 observation volumes; combined Model-MB and Model-T report R2 around 0.70 and relative RMSE below 20 percent, with strongest subvolume correlations up to R2 0.78. Pages pp.164-180; sections 2.1-2.8, 3, 4, 5; Tables 1-9. & M4; maturity; GT n; dedup no & Combining complementary UAV and ground RGB views with 3D point clouds improves coverage of tree subvolumes and supports tree-level flower-cluster estimation, illustrating the geometric route toward multi-view inventory aggregation. Limit: Small ground-truth sample and high biological variability; UAV point-cloud density falls toward lower tree sections; occlusion and sky-cloud confusion affect segmentation; no persistent cross-view identity or cluster-level deduplication; estimates flower intensity and clusters rather than a unique fruit inventory or maturity classes. & direct\\
Zhang, Wenli (2022) \cite{ledger_r_ee550a3e3a6fb318} & RGB\_\allowbreak{}video\_\allowbreak{}field\_\allowbreak{}rover\_\allowbreak{}DJI\_\allowbreak{}Osmo\_\allowbreak{}Action; citrus & Field rover at 2 m/s with 60 fps 1920x1080 video; 1465 cropped detection samples with 7:3 train-test split and 182 standard test images; OrangeYolo AP 0.938 and standard-image AP 0.957; six video sequences from two fields and 22 trees; OrangeSort region strategy MAE 0.081 and SD 0.08, compared with Sort MAE 0.45 and DeepSort MAE 1.212. Pages pp.1-14; sections Method, Experiment, Comparative analysis, Conclusion; Tables 1-5 and Figs.1-13. & M3; maturity; GT y; dedup explicit & Temporal tracking with a deliberate counting region and motion displacement estimation substantially reduces double counting in orchard video, providing direct evidence for the temporal component of unique fruit inventory construction. Limit: Tracking IDs are local to each video sequence rather than a persistent identity across sessions or trees; unseen canopy and fruit remain unobserved; performance degrades under poor lighting, overexposure, camera shake, and severe occlusion; no 3D spatial coordinates or maturity class; future work explicitly proposes 3D localization and unseen-fruit correction. & direct\\
Chen, Zhiyuan (2022) \cite{chen2022ffbdet} & RGB\_\allowbreak{}still\_\allowbreak{}images\_\allowbreak{}online\_\allowbreak{}dataset\_\allowbreak{}sources; oil palm FFB & Approximately 200 collected images with 169 suitable FFB images; deep-learning split 117 train, 32 validation, 20 test; YOLOv4 mAP 98.5 percent, precision 84 percent, recall 94 percent, F1 0.89, 31.8 ms; YOLOv5 mAP 96.4 percent, precision and recall 90.9 percent, F1 0.90, 11.4 ms. Pages pp.86-90; sections II-IV; Tables I-IV; Fig.2. & M0; maturity; GT n; dedup no & Comparative object detectors can localize multiple oil-palm FFBs in an image, with a speed-accuracy tradeoff between YOLOv4 and YOLOv5, but the study only establishes per-image detection. Limit: Images were sourced from Kaggle and Google rather than an actual plantation because of travel restrictions; small dataset and random image split; no field-disjoint evaluation, temporal tracking, cross-view identity, or deduplication; no maturity classes; authors identify real-plantation validation as future work. & direct\\
Phimpisan, Phaireepinas (2023) \cite{ledger_r_1c81e223679e4bfe} & RGB\_\allowbreak{}fluorescence\_\allowbreak{}images\_\allowbreak{}UV\_\allowbreak{}LED\_\allowbreak{}395\_\allowbreak{}to\_\allowbreak{}400\_\allowbreak{}nm; oil palm FFB & 100 individual FFB samples, 20 under-ripe and 80 ripe; five expert graders; 600 fruit samples with six fruits sampled per bunch; solvent-extraction CPO validation; reported 100 percent classification agreement and 2-5 second image capture. Pages pp.859-868; sections Materials and Methods, Results and Discussion, Conclusion; Tables 1-4; Figs.2-6. & M0; maturity; GT n; dedup no & Fluorescence imaging can classify the ripeness attribute of an individual oil-palm FFB rapidly and non-destructively, adding a sensor modality for the attribute layer of a class-wise inventory. Limit: Small and imbalanced sample with only two ripeness classes; controlled dark chamber or sorting-equipment setup; no field multi-view acquisition, detection of multiple bunches, temporal tracking, cross-view identity, or deduplication; the reported 100 percent result is not an external held-out field benchmark. & direct\\
Gené-Mola, Jordi (2023) \cite{ledger_r_2d972b94554c6cd1} & Kinect\_\allowbreak{}Azure\_\allowbreak{}RGB\_\allowbreak{}D\_\allowbreak{}or\_\allowbreak{}RGB\_\allowbreak{}video\_\allowbreak{}plus\_\allowbreak{}GNSS; apple & 420 trees in four rows, 84 five-meter stretches, two scanning dates, east and west sides, 1087 annotated frames; YOLOv5x F1 0.834 and mAP 0.871; ByteTrack MOTA 0.6817, IDF1 0.8369, HOTA 0.6894, 15.4 ms per image; yield mapping MAPE 8.45 percent and R2 0.66 at four weeks, MAPE 7.47 percent and R2 0.70 at one week. Pages pp.1-5; sections I-IV; Tables I-II; Figs.1-4. & M3; maturity; GT n; dedup explicit & Continuous motion video with MOT reduces duplicate counting across adjacent frames and supports stretch-level yield mapping, while appearance ReID did not improve tracking for visually similar apples. Limit: Track IDs are local to each scan and not persistent across separate dates or trees; ground truth is missing for several stretches; two-side counts are summed and calibrated rather than explicitly fused by 3D geometry; appearance-based ReID is weak for similar apples; no maturity/class attribute and no class-wise unique inventory. & direct\\
Shiddiq, Minarni (2023) \cite{shiddiq2023palmcount} & RGB\_\allowbreak{}video\_\allowbreak{}conveyor; oil palm FFB & 117 Tenera FFB, 10 videos, RGB 640x480 at 20 FPS, HSV saturation thresholding, contour area threshold above 9000, centroid association threshold below 50 pixels; per-video accuracy 50.00 to 100.00 percent, mean accuracy 79.35 percent, with recall, precision, and F-score reported. Pages pp.111-120; Abstract; Materials and Methods; Figs.1-6; Tables 1-2; Conclusion. & M3; maturity; GT n; dedup no & A low-cost RGB conveyor system can detect, track, and count moving oil palm FFB in real time, but performance is sensitive to illumination, visible FFB surface and orientation, adjacent bunches, and camera or conveyor timing. Limit: The experiment is a controlled postharvest conveyor setup rather than an in-canopy orchard view; one RGB camera and colour thresholding are sensitive to lighting and FFB orientation; track IDs are local to each video; there is no cross-view or cross-session identity, no tree-level unique inventory, and ripeness is explicitly excluded from evaluation. & direct\\
Rapado-Rincón, David (2023) \cite{ledger_r_700d34dae0f20d78} & RGB\_\allowbreak{}plus\_\allowbreak{}LiDAR\_\allowbreak{}structured\_\allowbreak{}point\_\allowbreak{}cloud\_\allowbreak{}multi\_\allowbreak{}view\_\allowbreak{}robot\_\allowbreak{}arm; tomato & Mask R-CNN training and evaluation uses 1180 images; world-model evaluation uses 700 images from 100 viewpoints of seven plants, 10 height levels, 960x540 colour images and structured point clouds; detection precision and recall are reported at IoU 0.5, along with MPE, MAPE, HOTA, DetA, and AssA; world-model update rate reaches 10 Hz; HOTA up to 71.47 percent and full-sequence MAPE 12.21 percent in the selected configuration. Pages pp.78-91; Abstract; Sections 2.2-2.5; Table 1; Table 3; Figs.2-8; Discussion; Conclusions. & M3/M4; maturity; GT y; dedup explicit & Multi-view perception lets later viewpoints recover tomatoes missed under occlusion, while 3D position association maintains an object-centric representation. The paper demonstrates that total count error alone can hide ID switches, false positives, and false negatives, so tracking metrics such as HOTA, DetA, and AssA are needed for a unique inventory representation. Limit: The evaluation covers seven greenhouse plants and a fixed semi-cylindrical viewpoint path rather than a broader orchard or field setting. Point-cloud noise and partial observations cause rejected detections and invalid sphere fits, especially in upper and occluded plant regions. Track IDs are evaluated within plant sequences, not as persistent identities across separate plants, dates, or sessions. The experiment uses one tomato class and does not evaluate maturity or class-wise inventory attributes. & direct\\
Abeyrathna, R. M. Rasika D. (2023) \cite{ledger_r_9e1316b934e66409} & RGB\_\allowbreak{}D\_\allowbreak{}Realsense\_\allowbreak{}D455\_\allowbreak{}plus\_\allowbreak{}video\_\allowbreak{}tractor; apple & 800 original images expanded to 6511 augmented images, split 4500/1500/511; 30 real apples attached to 10 artificial trees; three camera orientations at 0, 15, and 30 degrees and three speeds of 0.052, 0.069, and 0.098 m/s; model precision, recall, F1, and mAP@0.5; 3D distance reference measured for 30 apples three times; best RMSE 1.54 cm and counting accuracy up to 86.6 percent. Pages pp.1-19; Abstract; Sections 3.1-3.7; Table 1; Figs.8-20; Appendix A; Conclusion. & M3; maturity; GT n; dedup implicit & Combining RGB-D depth with deep detectors and Deep SORT supports dynamic apple localization and prevents the same apple from being counted repeatedly across consecutive frames. EfficientDet gives the lowest reported depth error at 15 degrees, while YOLOv5 and YOLOv7 provide the strongest counting result in the tested setup. Limit: The outdoor evaluation uses an artificially constructed orchard with only 30 apples on 10 artificial trees, so it does not establish performance in a natural orchard canopy. The experiment contains one apple class and no maturity or other class attribute, and track IDs are local to a moving sequence rather than persistent across sessions or trees. The augmented image split may not establish independent scene or tree-level test conditions, and the dataset is available only from the authors upon reasonable request. The study does not evaluate global cross-view identity or class-wise unique inventory. & direct\\
Hu, Jing (2023) \cite{ledger_r_d594d85dc382f9ff} & RGB\_\allowbreak{}images\_\allowbreak{}and\_\allowbreak{}video; apple & 4246 apple images from public datasets and free internet images, 80 percent training and 20 percent validation; YOLOv7, YOLOv7-CBAM, and YOLOv7-CA ablation; three test videos with 290 manual fruit counts in total; detection precision 0.924, recall 0.816, mAP@0.5 0.913, F1 0.85; proposed tracker relative MAE 0.064 overall. Pages pp.1-14; Abstract; Sections 2.1-2.5; Tables 1-3; Figs.1-9; Discussion; Conclusions. & M3; maturity; GT n; dedup implicit & Adding coordinate attention improves single-frame apple detection, while SURF appearance features combined with Kalman motion prediction and cascade matching reduces repeated counts across video frames compared with detector-only counting and SORT. Limit: The training data come from public datasets and free internet images, with pseudo-labeling and secondary annotation for some images; the tracking evaluation uses only three videos and does not report ID switch, MOTA, HOTA, or a full identity metric. Performance is unstable across lighting and shading conditions, and the paper does not establish tree-disjoint or field-disjoint generalization, 3D cross-view fusion, persistent identity across sessions, or class-wise inventory attributes. & direct\\
Lai, Jin Wern (2023) \cite{lai2023palmslr} & RGB\_\allowbreak{}colour\_\allowbreak{}vision\_\allowbreak{}LiDAR\_\allowbreak{}spectral\_\allowbreak{}NIR\_\allowbreak{}Raman\_\allowbreak{}fluorescence\_\allowbreak{}inductive\_\allowbreak{}thermal\_\allowbreak{}and\_\allowbreak{}other\_\allowbreak{}sensors; oil palm FFB & English-language Google Scholar search using oil-palm FFB maturity and detection terms, publication years 2012-2022; methods categorized by sample type, sensor, classifier, and deployment location; 11 unique approaches; abstract reports 51 papers and discussion reports 48 articles; colour vision appears in 22 articles; computer vision, LiDAR, and inductive sensors are identified as having plantation testing. Pages pp.1-16; Abstract; Sections 1-4; Tables 1-3; Figs.1-2; References. & M0; maturity; GT y; dedup n-a & Existing oil-palm review literature is centered on ripeness classification and grading, with colour vision plus CNN presented as a feasible field direction. It does not formulate the broader design space of multi-view identity, duplicate suppression, or class-wise unique fruit inventory. Limit: The search uses Google Scholar rather than a reproducible Scopus or Web of Science protocol and is restricted to English and 2012-2022. The paper is a secondary synthesis, not a primary benchmark, and contains an internal corpus-count discrepancy between the abstract and discussion. It does not evaluate cross-view association, persistent IDs, duplicate resolution, or tree-level inventory construction. & direct\\
Jordan A. James (2023) \cite{ledger_r_fd5751e2ca2f7496} & RGB\_\allowbreak{}still\_\allowbreak{}images\_\allowbreak{}ground\_\allowbreak{}based\_\allowbreak{}tree\_\allowbreak{}views; citrus & 579 images; 44,233 full-dataset bounding boxes; 60 varieties and 8 species/genera; 187 trees; tiled YOLOv7 AP 45.5, AP50 83.1; tree-yield front+back R2 0.754 detect-count and 0.793 filter-detect-count Pages 1-7. & M0; maturity; GT n; dedup no & A public citrus benchmark makes on-tree versus fallen fruit explicit and improves tree-yield correlation through canopy filtering, but front/back aggregation does not establish unique fruit identity. Limit: No cross-view fruit identity, tracking, ID-switch, HOTA, or ReID evaluation; front/back counts are summed by assumption; fully occluded fruit are unlabeled; small-object errors; no maturity classes; random train/test image split and tree-yield correlation are not equivalent to unique inventory. & direct\\
Lin, Yishen (2024) \cite{ledger_r_1c69e46dac45d45d} & RGB\_\allowbreak{}still\_\allowbreak{}images\_\allowbreak{}ground\_\allowbreak{}based\_\allowbreak{}orchard; citrus & 957 images; 670/190/97 train-validation-test; AG-YOLO P 90.6\%, R 73.4\%, mAP50 83.2\%, mAP50:95 60.3\%, APs 18.8\%, APm 47.2\%, APl 69.2\%, 34.22 FPS Pages 1,4-5,9-10,13-14. & M0; maturity; GT y; dedup no & Global context fusion with NextViT improves single-view citrus detection under severe occlusion, but does not resolve cross-view identity or duplicate counting. Limit: Single-view RGB benchmark; no tracking, ReID, duplicate resolution, or unique-inventory metric; raw data only on request; ripe/unripe is not identity; small and selected dataset. & direct\\
Ambrus, B. (2024) \cite{ledger_r_3e129b7285059b95} & RGB\_\allowbreak{}DSLR\_\allowbreak{}and\_\allowbreak{}robot\_\allowbreak{}RGB\_\allowbreak{}plus\_\allowbreak{}LiDAR\_\allowbreak{}calibration\_\allowbreak{}plus\_\allowbreak{}Ciclop\_\allowbreak{}3D\_\allowbreak{}scan; tomato & Robot validation F1 59.3\%; red F1 0.70 and green F1 0.22; 3D model weight relative error 21.90\% with SD 17.9665\%; DSLR all-treatment weight R2 0.867; robot all-treatment weight R2 0.907; DSLR all-treatment count R2 0.744; robot all-treatment count R2 0.806 Pages 1,4-6,9-12,14-16,18. & M0/M4; maturity; GT n; dedup no & Combining robot/DSLR RGB perception with 3D fruit shape supports field tomato count and weight estimation, but canopy occlusion and one-sided views cause missing fruit and prevent unique inventory. Limit: No persistent fruit ID, cross-view ReID, duplicate resolution, ID-switch/HOTA, or class-wise global inventory; small and treatment-specific field experiment; weak green-tomato classification; one-sided robot view and canopy occlusion; data available on request; 3D scanner models harvested fruit shape rather than associating views. & direct\\
Vélez, Sergio (2024) \cite{ledger_r_659a7da54b8a43e6} & geotagged\_\allowbreak{}smartphone\_\allowbreak{}RGB\_\allowbreak{}plus\_\allowbreak{}UAV\_\allowbreak{}multispectral\_\allowbreak{}orthomosaic\_\allowbreak{}and\_\allowbreak{}SfM\_\allowbreak{}point\_\allowbreak{}cloud\_\allowbreak{}plus\_\allowbreak{}RTK\_\allowbreak{}GNSS; grape & 253 smartphone images; 10 UAV 3D point clouds; 2 orthomosaics; RTK trunk locations; per-plant cluster counts for 42 plants; open Zenodo dataset Pages 1-5,7-10,13. & M4; maturity; GT n; dedup n-a & Geospatial plant anchors can align ground imagery, aerial multimodal data, and plant-level fruit-cluster counts, providing an agricultural transfer pattern for cross-observation registration. Limit: Identity is at plant level, not individual fruit or bunch level; no fruit-level detection/tracking/ReID or duplicate-resolution evaluation; dataset article does not benchmark a unique inventory algorithm; two vineyards and one vine cultivar in one Spanish region; raw archives are large. & transferable\\
Wu, Haitao (2024) \cite{ledger_r_92ac7cb4b797197d} & RGB\_\allowbreak{}still\_\allowbreak{}images\_\allowbreak{}natural\_\allowbreak{}orchard\_\allowbreak{}plus\_\allowbreak{}synthetic\_\allowbreak{}fog; apple & 10,304 images; 8,240/1,032/1,032 train-validation-test; overall P 90.7\%, R 88.9\%, mAP50 94.3\%, 10.46M parameters, 25.4 GFLOPs, 2.9 ms; per-weather mAP50 95.2/95.4/95.3/88.9 Pages official full-text HTML: abstract, sections 2.1.3, 3, 3.4-3.6, conclusion. & M0; maturity; GT n; dedup no & A lightweight detector remains accurate across several weather conditions, but fog and single-angle occlusion remain unresolved and the result is still intra-view detection. Limit: No tracking, ReID, cross-view identity, duplicate resolution, ID-switch/HOTA, or unique inventory metric; one apple variety; distant fruit unlabeled; synthetic fog augmentation and single-angle limitations; no class-wise maturity attribute. & direct\\
N P Vidhya (2024) \cite{ledger_r_e8d88a31640ee300} & RGB\_\allowbreak{}mobile\_\allowbreak{}camera\_\allowbreak{}still\_\allowbreak{}images; mango & 473 images; random 80/10/10 split; YOLOv7; image 640x640; batch 16; 100 epochs; reported mAP@0.5 99.5\% and mAP@0.5:0.95 81.9\% in training/validation phase Pages 1,3-6. & M0; maturity; GT y; dedup no & A simple transfer-learned YOLOv7 detector gives high reported single-view mango detection metrics, but evidence does not address cross-view unique counting. Limit: No tracking, ReID, cross-view identity, duplicate resolution, ID-switch/HOTA, or unique inventory; random split; healthy mangoes only; one class; no separate numerical test metrics; reported accuracy context is training/validation. & direct\\
Hobart, Marius (2025) \cite{ledger_m_66042c905cc191de} & UAV\_\allowbreak{}RGB\_\allowbreak{}multiview\_\allowbreak{}photogrammetric\_\allowbreak{}dense\_\allowbreak{}point\_\allowbreak{}cloud\_\allowbreak{}plus\_\allowbreak{}GNSS\_\allowbreak{}RTK; apple & 0.5 ha; 1334 trees; 2 seasons; 4 rows per season; 14 blocks of 8 trees; 7.5 m and 10 m flights; apple count R2 about 0.41, 0.40, 0.53; main yield-mass R2 0.56 at 7.5 m and 0.76 at 10 m; one 7.5 m 2020 cloud failed Pages 1-3,5-8,10-16. & M4; maturity; GT n; dedup no & Low-altitude UAV RGB SfM can recover ripe-apple 3D positions and volumes for block-level yield-mass mapping, with 10 m outperforming 7.5 m in the main mass result. Limit: No persistent fruit identity, cross-view ReID, duplicate resolution, ID-switch/HOTA, or class-wise unique inventory; block-level validation; leaf occlusion and noisy/shifted point clouds; one failed cloud and excluded outer rows; ripe red fruit only; reported conclusion R2 is inconsistent with the main results; data on request. & direct\\
Guillem Montalbán Faet (2025) \cite{ledger_m_949f38a7336860d8} & UAV\_\allowbreak{}RGB\_\allowbreak{}plus\_\allowbreak{}G\_\allowbreak{}RE\_\allowbreak{}NIR\_\allowbreak{}multispectral\_\allowbreak{}nadir\_\allowbreak{}images; citrus & 550 synchronized captures; 2,750 images; 301,232 annotations; 14 m AGL; GSD 0.65 cm/pixel; 80/70 overlap; four annotators plus two reviewers; date-based 2-train/1-test recommendation; 70/15/15 alternative Pages 1-4,6-10. & M4; maturity; GT n; dedup no & A public UAV dataset makes cross-spectral fruit annotations geometrically reproducible and provides a date-based evaluation protocol that can reduce leakage. Limit: Data article without detector benchmark; no persistent fruit identity across dates/views, ReID, duplicate resolution, ID-switch/HOTA, or unique inventory; one orchard and altitude; controlled midday illumination; large storage and SfM/reprojection cost; no multi-altitude evaluation. & direct\\
Soham Dutta (2025) \cite{ledger_m_dcfc5d3cb912146d} & RGB\_\allowbreak{}UAV\_\allowbreak{}ESP32\_\allowbreak{}CAM\_\allowbreak{}plus\_\allowbreak{}RaspberryPi\_\allowbreak{}edge\_\allowbreak{}inference; apple & Three collections about 10,000 samples each; apple detection split 79/21; YOLOv8 300 epochs; validation P 86.1\%, R 85.3\%, F1 85.7\%, mAP50 91\%, mAP50:95 63.1\%; no numerical YOLO test metrics Pages 1-7. & M0; none; GT y; dedup no & A low-cost RGB edge pipeline integrates orchard perception tasks, but it does not yet implement temporal consistency required for unique fruit inventory. Limit: Preprint; incomplete dataset provenance and field protocol; no temporal tracking, ReID, cross-view identity, duplicate resolution, ID-switch/HOTA, or unique inventory; single fruit class; detection test metrics absent; temporal consistency is future work. & direct\\
Ainul Hakim Fizal Sabillah (2025) \cite{ledger_r_134385f2faed38a6} & RGB\_\allowbreak{}smartphone\_\allowbreak{}and\_\allowbreak{}UAV\_\allowbreak{}video; oil palm FFB & 119 images; about 8 minutes video; about 22 trees; YOLOv8 416 precision 0.96, recall 0.97, F1 0.97, mAP 98.98\%; YOLOv7 416 mAP 98.87\%; augmentation YOLOv3 mAP 71.16\% vs 32.54\% without Pages 1-6,9-16. & M0; maturity; GT y; dedup no & Input resolution and augmentation strongly affect FFB detection, and newer YOLO versions report higher single-view metrics, but the evidence does not establish unique bunch inventory. Limit: Small dataset; no clear independent test/tree-disjoint split; best-iteration metrics may be training/validation; one FFB class and no ripeness labels; no tracking, ReID, duplicate resolution, ID-switch/HOTA, or unique inventory; circular video not used for identity. & direct\\
Goh, Jin Yu (2025) \cite{goh2025palmdataset} & RGB\_\allowbreak{}plus\_\allowbreak{}RealSense\_\allowbreak{}depth\_\allowbreak{}plus\_\allowbreak{}point\_\allowbreak{}cloud; oil palm FFB & 400 RGB images; 400 depth maps; 400 point clouds; four Johor locations; registration mean error 1.8 cm at 3 m; DBSCAN cluster validation; table agreement 87.5\% initial, 100\% consensus, 95.2\% recheck; abstract reports 92.5\% Pages 1-9. & M4; maturity; GT n; dedup implicit & A natural-plantation RGB-depth-point-cloud resource supports instance-level FFB ripeness and 3D localization with explicit calibration, a strong foundation for multimodal inventory research. Limit: Data descriptor without detector benchmark; no temporal/multi-view sequence identity or global deduplication; scene-local cluster IDs; 400 scenes only; tripod-style acquisition and no cross-session inventory; internal agreement figures differ between abstract and validation table. & direct\\
Timprae, Warut (2025) \cite{ledger_r_4d09eb83eea41b2f} & RGB\_\allowbreak{}handheld\_\allowbreak{}video\_\allowbreak{}plus\_\allowbreak{}SfM\_\allowbreak{}MVS\_\allowbreak{}Nerfacto\_\allowbreak{}3D\_\allowbreak{}point\_\allowbreak{}cloud; tomato & Laboro Tomato 804 images: YOLOv8x-Seg P=0.806 R=0.818 F1=0.812 and table segment mAP@0.5=0.881; text inconsistently states mAP at IoU 0.8. 3D diameter average error=8.01\% versus caliper; summer size gmax=1.08 versus winter=0.87 and ripeness gmax=0.41 versus 0.21. Pages 1-2,4-5,7-10,12-20. & M4; maturity; GT n; dedup implicit & Combines per-fruit instance segmentation, SfM/MVS/Nerfacto reconstruction, DBSCAN grouping, ellipsoid fitting, metric reference-ball calibration and CIELAB ripeness analysis for longitudinal greenhouse phenotyping. Limit: One greenhouse setting and cherry-tomato cultivar; handheld uncalibrated camera with about 180-degree view; occlusion and lighting affect reconstruction; object selection, shape verification and tracking still require manual intervention; no global fruit re-identification, ID-switch, HOTA or cross-view deduplication metrics; data available upon request. & direct\\
N. A. Hamid (2025) \cite{ledger_r_504a1d386928344f} & RGB\_\allowbreak{}computer\_\allowbreak{}vision\_\allowbreak{}CNN\_\allowbreak{}YOLO\_\allowbreak{}MSI\_\allowbreak{}HSI\_\allowbreak{}physical\_\allowbreak{}sensors\_\allowbreak{}IoT\_\allowbreak{}and\_\allowbreak{}edge\_\allowbreak{}computing; oil palm FFB & Secondary comparison across 46 references with summary tables for computer vision, YOLO, MSI/HSI and sensor systems; reported ranges include YOLO 91-96.7\%, MSI/HSI 90-95\%, and sensors 88-94.2\%. No systematic database query, screening flow, dataset-level reanalysis or independent pooled evaluation is reported. Pages 1-9,10-12. & M0; maturity; GT n; dedup n-a & Prior oil-palm review work concentrates on ripeness grading, sensors, spectral imaging, CNNs and YOLO deployment; it does not formulate cross-view identity, duplicate suppression or class-wise unique fruit inventory as a primary review problem. Limit: Narrative mini-review with no reproducible search protocol or explicit inclusion/exclusion process; reported metrics are inherited from cited studies and cannot be merged with the current corpus; no cross-view association, persistent identity, duplicate resolution or tree-level unique inventory analysis; some cited claims require verification against primary studies. & direct\\
Minarni Shiddiq (2025) \cite{ledger_r_722bfce242f37ecb} & RGB\_\allowbreak{}CMOS\_\allowbreak{}camera\_\allowbreak{}ImageJ\_\allowbreak{}threshold\_\allowbreak{}segmentation\_\allowbreak{}and\_\allowbreak{}particle\_\allowbreak{}analysis; oil palm FFB & 40 FFBs: 20 Dura and 20 Tenera, each with 10 unripe and 10 ripe; 80 front/back images; RGB intensity and ImageJ outer-fruit counts compared with manual counts. Mass-density regression R2: Dura ROI 0.4038, Tenera ROI 0.0384, Dura segmentation 0.2803, Tenera segmentation 0.0888. Pages 1-9. & M0; maturity; GT n; dedup no & Tenera shows higher RGB intensity, outer-fruit count and fruitlet density than Dura; front/back ImageJ counts are lower than manual counts because overlapping fruitlets are not separated reliably. Limit: Small sample from one smallholder plantation; controlled white conveyor background and fixed 60 cm camera; only outer fruits visible in each side view; no cross-view registration or unique fruit inventory; ImageJ thresholding and manual ROI require intervention; authors note overlapping fruits need watershed or more advanced processing; no classifier benchmark. & direct\\
Lee, Jaehwan (2026) \cite{ledger_r_d4270ea8ab11f268} & RGB\_\allowbreak{}LiDAR\_\allowbreak{}plus\_\allowbreak{}FAST\_\allowbreak{}LIO2\_\allowbreak{}SLAM; pear & Training, validation, and test sets contain 9000, 2580, and 100 images with 543375 labels; field validation on 1891 fruits reports 96.2 percent counting accuracy, 3.8 percent relative error, and MAE 0.64 fruits per square metre. Pages 2,5,8-10. & M3/M4; none; GT n; dedup explicit & Instance-level RGB and LiDAR fusion combines 2D detection with SLAM point clouds and temporal association to reduce duplicate counts while mapping fruit density. Limit: Single pear class and one horizontal-trellis orchard; 53 double-counts remain from association threshold and SLAM drift, so transfer to dense oil-palm canopy and class-wise inventory is not demonstrated. & direct\\
Meyer, Lukas (2024) \cite{ledger_m_1210e3ea00486b68} & RGB\_\allowbreak{}multiview\_\allowbreak{}plus\_\allowbreak{}SfM\_\allowbreak{}semantic\_\allowbreak{}NeRF; agricultural fruit or transferable mechanism & Synthetic six-fruit dataset: F1 0.95 with ground-truth masks and about 0.88 with SAM masks; real apple trees above 89 percent detection and Fuji benchmark F1 0.79. Pages 1-2,4-8. & M4; none; GT n; dedup explicit & Semantic NeRF converts unordered posed RGB observations into a 3D fruit field and clustered fruit point cloud, directly addressing duplicate counting across views. Limit: Manual per-tree cropping is required; clustering is sensitive to hyperparameters; performance drops under sparse views and the study does not evaluate class-wise maturity inventory. & direct\\
Wang, Penggang (2025) \cite{ledger_r_b2e0e9aabad17a97} & RGB\_\allowbreak{}D\_\allowbreak{}plus\_\allowbreak{}ORB\_\allowbreak{}SLAM3\_\allowbreak{}BiSeNetV2\_\allowbreak{}OctoMap; tomato & 600 annotated tomato images; fruit counts on three plants and three rows report mean relative error 14.36 percent; fruit volume on one plant reports mean relative error 14.25 percent; semantic segmentation mIoU 95.37 percent and field mATE 0.16 m. Pages 5-7,12,16-18,21. & M4; size; GT n; dedup implicit & RGB-D semantic SLAM projects fruit segmentation into a 3D point cloud, filters noise, and uses spherical fitting to estimate per-plant and per-row fruit counts and fruit volume. Limit: The study is not a cross-view deduplication evaluation; leaf and fruit occlusion plus depth errors cause missed or distorted fruit geometry, and omitted fitted fruits are excluded from volume statistics. & direct\\
\end{longtable}
\makeatletter
\@ifclassloaded{IEEEtran}{\twocolumn}{ }
\makeatother


\section{Exact search strings}
\label{app:queries2}

The following blocks reproduce the final Scopus strings. All seven use \texttt{PUBYEAR > 2014 AND PUBYEAR < 2027} and the subject-area clause shown below. The line breaks are for typesetting only; the tokens are unchanged.

\subsection*{Q1: Multi-view fruit inventory}
\begin{verbatim}
TITLE-ABS-KEY(("fruit" OR "fresh fruit bunch" OR "FFB" OR "oil palm"
OR "apple" OR "citrus" OR "mango" OR "grape" OR "berry" OR "bunch"
OR "cluster" OR "crop")
AND ("count*" OR "yield estimation" OR "load estimation"
OR "fruit load" OR "inventory" OR "enumeration")
AND ("multi-view" OR "multiple view*" OR "multi-camera" OR "video"
OR "cross-view" OR "structure from motion" OR "SfM" OR "track*"
OR "3D reconstruction" OR "image sequence"))
AND PUBYEAR > 2014 AND PUBYEAR < 2027
AND (SUBJAREA(AGRI) OR SUBJAREA(COMP) OR SUBJAREA(ENGI)
OR SUBJAREA(MULT))
\end{verbatim}

\subsection*{Q2: Cross-view identity and duplicates}
\begin{verbatim}
TITLE-ABS-KEY(("multi-view" OR "cross-view" OR "multi-camera"
OR "multiple cameras" OR "overlapping view*" OR "multi-sensor")
AND ("re-identification" OR "data association" OR "identity"
OR "duplicate*" OR "deduplicat*" OR "double counting"
OR "correspondence" OR "instance matching" OR "track*"
OR "association")
AND ("object detection" OR "instance segmentation" OR "counting"
OR "instance" OR "pedestrian" OR "vehicle"))
AND PUBYEAR > 2014 AND PUBYEAR < 2027
AND (SUBJAREA(COMP) OR SUBJAREA(ENGI) OR SUBJAREA(AGRI)
OR SUBJAREA(MULT))
\end{verbatim}

\subsection*{Q3: Plant geometry and multimodal sensing}
\begin{verbatim}
TITLE-ABS-KEY(("RGB-D" OR "depth camera" OR "stereo vision" OR "LiDAR"
OR "point cloud" OR "monocular depth" OR "photogrammetry"
OR "structure from motion")
AND ("fruit" OR "orchard" OR "vineyard" OR "canopy" OR "plant"
OR "tree crop" OR "plantation" OR "oil palm")
AND ("detection" OR "segmentation" OR "localization" OR "counting"
OR "phenotyping" OR "harvesting"))
AND PUBYEAR > 2014 AND PUBYEAR < 2027
AND (SUBJAREA(AGRI) OR SUBJAREA(COMP) OR SUBJAREA(ENGI)
OR SUBJAREA(EART) OR SUBJAREA(ENVI) OR SUBJAREA(MULT))
\end{verbatim}

\subsection*{Q4: Per-instance attributes}
\begin{verbatim}
TITLE-ABS-KEY(("fruit" OR "produce" OR "crop" OR "bunch")
AND ("maturity" OR "ripeness" OR "grading" OR "grade"
OR "quality class*" OR "size class*" OR "defect" OR "disease"
OR "cultivar" OR "variety")
AND ("instance segmentation" OR "object detection" OR "per-instance"
OR "individual fruit" OR "bounding box" OR "per-object"))
AND PUBYEAR > 2014 AND PUBYEAR < 2027
AND (SUBJAREA(AGRI) OR SUBJAREA(COMP) OR SUBJAREA(ENGI)
OR SUBJAREA(MULT))
\end{verbatim}

\subsection*{Q5: Previous reviews}
\begin{verbatim}
TITLE(("review" OR "survey" OR "systematic review" OR "scoping review"))
AND TITLE-ABS-KEY(("fruit detection" OR "fruit counting"
OR "yield estimation" OR "oil palm" OR "orchard"
OR "crop monitoring" OR "agricultural robot*"
OR "precision agriculture")
AND ("deep learning" OR "computer vision" OR "object detection"
OR "machine vision"))
AND PUBYEAR > 2014 AND PUBYEAR < 2027
AND (SUBJAREA(AGRI) OR SUBJAREA(COMP) OR SUBJAREA(ENGI)
OR SUBJAREA(MULT))
\end{verbatim}

\subsection*{Q6: Oil-palm vision and machine learning}
\begin{verbatim}
TITLE-ABS-KEY(("oil palm" OR "elaeis guineensis"
OR "fresh fruit bunch" OR "FFB")
AND ("deep learning" OR "machine learning" OR "convolutional"
OR "neural network" OR "computer vision" OR "image processing"
OR "YOLO" OR "object detection" OR "remote sensing"))
AND PUBYEAR > 2014 AND PUBYEAR < 2027
AND (SUBJAREA(AGRI) OR SUBJAREA(COMP) OR SUBJAREA(ENGI)
OR SUBJAREA(EART) OR SUBJAREA(ENVI) OR SUBJAREA(MULT))
\end{verbatim}

\subsection*{Q7: Single-view fruit load}
\begin{verbatim}
TITLE-ABS-KEY(("fruit" OR "apple" OR "citrus" OR "mango" OR "grape"
OR "berry" OR "bunch" OR "oil palm" OR "orchard" OR "vineyard")
AND ("counting" OR "count" OR "yield estimation" OR "load estimation"
OR "fruit load" OR "crop load")
AND ("deep learning" OR "convolutional" OR "object detection"
OR "YOLO" OR "instance segmentation" OR "machine vision"))
AND PUBYEAR > 2014 AND PUBYEAR < 2027
AND (SUBJAREA(AGRI) OR SUBJAREA(COMP) OR SUBJAREA(ENGI)
OR SUBJAREA(MULT))
\end{verbatim}

\bibliographystyle{IEEEtran}
\bibliography{references3}
\end{document}
